\documentclass[t, aspectratio=169]{beamer}
\usepackage{fontspec}
\setmainfont{Times New Roman}
\usepackage{amsmath}
\usepackage{amssymb}
\usepackage{booktabs}
\usepackage{graphicx}
\usepackage{xcolor}
\usepackage{array}
\usepackage{tikz}
\usetikzlibrary{shapes.callouts, positioning, arrows.meta, fit, backgrounds}

\definecolor{hustblue}{RGB}{0,60,113}
\definecolor{goodgreen}{RGB}{30,120,60}
\definecolor{warnorange}{RGB}{200,90,20}
\definecolor{badred}{RGB}{170,30,30}

\def\slideContentLeftMargin{0.5cm}
\def\slideContentRightMargin{0.5cm}
\def\hustHeaderHeight{3.0cm}
\def\slideContentBotMargin{2.8cm}

\usepackage[blue, 169]{beamerthemeHUST}
\usefonttheme{serif}
\renewcommand{\familydefault}{\rmdefault}

\setbeamercolor{frametitle}{fg=hustblue}
\setbeamerfont{frametitle}{size=\fontsize{16}{18}\selectfont, series=\bfseries}
\setbeamertemplate{frametitle}{
  \nointerlineskip
  \begin{tikzpicture}[remember picture, overlay]
    \node[anchor=north west, xshift=0.35cm, yshift=-0.95cm, text width=12.3cm, align=left]
    at (current page.north west)
    {\usebeamerfont{frametitle}\usebeamercolor[fg]{frametitle}\insertframetitle};
  \end{tikzpicture}
}
\setbeamertemplate{footline}{\vbox to 1.8cm{}}

\renewcommand{\titlePageTitleX}{-0.5cm}
\renewcommand{\titlePageTitleY}{-0.5cm}
\renewcommand{\hustTitleAlign}{\alignLeft}
\renewcommand{\hustSubtitleAlign}{\alignLeft}
\renewcommand{\hustAuthorAlign}{\alignLeft}
\renewcommand{\hustInstAlign}{\alignLeft}

\renewcommand{\hustsectionpage}[1]{%
  {%
    \setbeamertemplate{background}{\includegraphics[width=\paperwidth,height=\paperheight]{\hustbg{section}}}%
    \begin{frame}[plain,c]
      \centering \Huge \textcolor{white}{\textbf{#1}}
    \end{frame}
  }%
}

\graphicspath{{images/}}

% ---- Assertion-Evidence helper: big-number card, direct-labeled ----
\newcommand{\bignum}[4]{%
  % #1 = number/value (big), #2 = short caption below, #3 = color, #4 = width
  \begin{minipage}[t]{#4}
    \centering
    {\fontsize{32}{36}\selectfont\bfseries\color{#3} #1}\\[2pt]
    {\small #2}
  \end{minipage}%
}
% Compact variant for slides that also carry a table/diagram below the numbers
\newcommand{\bignumS}[4]{%
  \begin{minipage}[t]{#4}
    \centering
    {\fontsize{22}{26}\selectfont\bfseries\color{#3} #1}\\[1pt]
    {\scriptsize #2}
  \end{minipage}%
}

\title{\large Tích hợp dữ liệu khí tượng đa nguồn}
\subtitle{\normalsize Phát hiện dị thường khí hậu \& hỗ trợ ra quyết định đa ngành khu vực Đông Nam Á}
\author{Ngô Văn Quang \quad $\cdot$ \quad Nguyễn Cao Lương \quad $\cdot$ \quad Trương Tuấn Nghĩa}
\institute{SoICT, Đại học Bách khoa Hà Nội \\
IT5427 -- Tích hợp và xử lý dữ liệu lớn -- Nhóm 16\\[3pt]
{\small Giảng viên: TS.\ Vũ Tuyết Trinh}}

\begin{document}

% 1. Title
\husttitlepage

% 2. Outline
\begin{frame}[shrink=5]{Nội dung}
\begin{columns}[T]
\begin{column}{0.5\textwidth}
\begin{enumerate}
\item \hyperlink{sec1}{Vấn đề \& mục tiêu đề tài}
\item \hyperlink{sec2}{Nguồn dữ liệu -- đã kiểm chứng thật}
\item \hyperlink{sec3}{Kiến trúc hệ thống \& tích hợp}
\end{enumerate}
\end{column}
\begin{column}{0.5\textwidth}
\begin{enumerate}
\setcounter{enumi}{3}
\item \hyperlink{sec4}{Trạm quan trắc \& dung lượng dữ liệu}
\item \hyperlink{sec5}{Kế hoạch \& đội hình}
\item \hyperlink{sec6}{Rủi ro \& kết luận}
\end{enumerate}
\end{column}
\end{columns}
\end{frame}

\resetPageCount
\showPageNumber
\useStyleHustFull
\useThinBar

%=============================================================
\section{Vấn đề \& mục tiêu}
%=============================================================
\hypertarget{sec1}{}
\hustsectionpage{1. Vấn đề \& mục tiêu\\đề tài}

% Frame 3
\begin{frame}{Dữ liệu khí tượng đa nguồn buộc phải giải quyết đúng 2 loại dị biệt cốt lõi của môn học}
\begin{center}
\resizebox{0.86\textwidth}{!}{%
\begin{tikzpicture}[remember picture,
  box/.style={rectangle, rounded corners, draw=hustblue, thick, fill=hustblue!8, align=center, inner sep=6pt, font=\small},
  arr/.style={-{Latex[length=3mm]}, thick, hustblue}]
\node[box, minimum width=5.5cm, minimum height=1.6cm] (p1) at (-4.2,1.6) {\textbf{Structural/Semantic Heterogeneity}\\5 nguồn: lưới liên tục (ERA5) vs.\\điểm rời rạc (trạm), field/đơn vị khác nhau};
\node[box, minimum width=5.5cm, minimum height=1.6cm, fill=warnorange!10, draw=warnorange] (p2) at (4.2,1.6) {\textbf{Instance/Identity Heterogeneity}\\Cùng 1 đại lượng (nhiệt độ, mưa...)\\được đo bởi tới 5 nguồn trùng nhau};
\node[box, minimum width=6.2cm, minimum height=1.3cm, fill=goodgreen!10, draw=goodgreen] (s1) at (-4.2,-1.3) {Giải pháp: \textbf{Mediated Schema (GAV)}\\+ Spatial Interpolation};
\node[box, minimum width=6.2cm, minimum height=1.3cm, fill=goodgreen!10, draw=goodgreen] (s2) at (4.2,-1.3) {Giải pháp: \textbf{Survivorship Rule}\\(Record Linkage / Data Fusion)};
\draw[arr] (p1) -- (s1);
\draw[arr] (p2) -- (s2);
\end{tikzpicture}%
}
\end{center}
\end{frame}

% Frame 4
\begin{frame}{Đề tài phục vụ 3 nhóm ứng dụng cụ thể, không dừng lại ở việc gộp dữ liệu}
\centering
\begin{tabular}{@{}p{3.6cm}p{8.6cm}@{}}
\toprule
Ứng dụng & Cách dùng kho dữ liệu đã chuẩn hoá \\
\midrule
Phát hiện dị thường khí hậu & So sánh giá trị hiện tại với phân bố lịch sử cùng khu vực/mùa vụ \\
Nông nghiệp \& năng lượng & Ước lượng tiềm năng bức xạ mặt trời/gió theo khu vực và thời điểm \\
Vận hành UAV an toàn & Khuyến nghị khung giờ/khu vực bay dựa trên gió và lượng mưa dự kiến \\
\bottomrule
\end{tabular}
\vspace{0.3cm}

\small Cả 3 ứng dụng dùng chung một \textbf{kho dữ liệu khí tượng đã chuẩn hoá} (Mediated Schema) --
xây một lần, tái sử dụng nhiều lần, đúng tinh thần Data Integration của môn học.
\end{frame}

%=============================================================
\section{Nguồn dữ liệu}
%=============================================================
\hypertarget{sec2}{}
\hustsectionpage{2. Nguồn dữ liệu --\\đã kiểm chứng thật}

% Frame 5
\begin{frame}{8/9 nguồn dữ liệu đã gọi API/tải mẫu thật thành công, không chỉ tra tài liệu}
\begin{center}
\bignumS{8/9}{nguồn kiểm tra thành công\\(script \texttt{check\_data\_sources.py}, chạy thật 30/9/2026)}{goodgreen}{5.6cm}
\hspace{0.5cm}
\bignumS{4/4}{nguồn lõi hoạt động ngay\\KHÔNG cần đăng ký tài khoản}{goodgreen}{5.6cm}
\end{center}
\vspace{0.15cm}
\centering\scriptsize
\begin{tabular}{@{}llll@{}}
\toprule
Nguồn & Trạng thái & Cần đăng ký? & Bằng chứng \\
\midrule
NASA POWER & \textcolor{goodgreen}{\textbf{OK}} & Không & HTTP 200, 5 ngày dữ liệu Hà Nội \\
Open-Meteo Historical & \textcolor{goodgreen}{\textbf{OK}} & Không & HTTP 200, 5 ngày dữ liệu Jakarta \\
NOAA GHCN-Daily (S3) & \textcolor{goodgreen}{\textbf{OK}} & Không & 44.560 bản ghi, 1958--2025 \\
NOAA ISD/Global Hourly & \textcolor{goodgreen}{\textbf{OK}} & Không & 16.741 bản ghi/năm, 1 trạm \\
ERA5 (ARCO, Google Cloud) & \textcolor{goodgreen}{\textbf{OK}} & Không & 35 biến, \texttt{.zmetadata} tải được \\
\bottomrule
\end{tabular}
\end{frame}

% Frame 6
\begin{frame}{Kiểm tra thật đã sửa 3 lỗi và phát hiện 2 nguồn ERA5 miễn phí bị bỏ sót}
\centering\footnotesize
\begin{tabular}{@{}p{3.3cm}p{9.3cm}@{}}
\toprule
Phát hiện & Nội dung \\
\midrule
\textcolor{badred}{\textbf{Sửa}} ERA5 trên AWS & \texttt{era5-pds} là bucket \textit{Requester Pays} -- cần tài khoản AWS thật, KHÔNG miễn phí hoàn toàn như ghi chú ban đầu \\
\textcolor{badred}{\textbf{Sửa}} OpenAQ & API v1/v2 đã bị khai tử (HTTP 410 Gone) -- phải dùng v3, key qua header khác \\
\textcolor{goodgreen}{\textbf{Đính chính}} Himawari-8 & Vẫn hoạt động thật (đã tải mẫu file 12/10/2025) -- KHÔNG ngừng vận hành như nhận định sai trước đó \\
\textcolor{hustblue}{\textbf{Mới}} ARCO-ERA5 & Google Cloud, hoàn toàn miễn phí, không đăng ký -- 35 biến, 1959--2022, 0.25\textdegree{}, hourly \\
\textcolor{hustblue}{\textbf{Mới}} Planetary Computer & Microsoft, cấp SAS token tự động, không đăng ký tài khoản \\
\bottomrule
\end{tabular}
\end{frame}

%=============================================================
\section{Kiến trúc \& tích hợp}
%=============================================================
\hypertarget{sec3}{}
\hustsectionpage{3. Kiến trúc hệ thống\\\& tích hợp dữ liệu}

% Frame 7
\begin{frame}{Pipeline xử lý 5 nguồn hội tụ về một Mediated Schema duy nhất trước khi phân tích}
\begin{center}
\resizebox{0.92\textwidth}{!}{%
\begin{tikzpicture}[remember picture,
  f/.style={rectangle, rounded corners, draw=hustblue, thick, fill=hustblue!8, minimum width=2.6cm, minimum height=0.9cm, align=center, font=\scriptsize},
  arr/.style={-{Latex[length=2.5mm]}, thick, hustblue}]
\node[f] (s1) at (-6.5,2.2) {NASA POWER};
\node[f] (s2) at (-6.5,1.0) {Open-Meteo};
\node[f] (s3) at (-6.5,-0.2) {GHCN-Daily\\(PIVOT)};
\node[f] (s4) at (-6.5,-1.4) {ISD\\(PARSE)};
\node[f] (s5) at (-6.5,-2.6) {ERA5 (ARCO)};
\node[f, minimum width=3.2cm, minimum height=2.6cm, fill=warnorange!10, draw=warnorange] (raw) at (-2.2,-0.2) {Raw Zone\\+ chuẩn hoá\\đơn vị/thời gian};
\node[f, minimum width=3.2cm, minimum height=2.6cm, fill=goodgreen!10, draw=goodgreen] (ms) at (1.6,-0.2) {Mediated\\Schema\\(Survivorship)};
\node[f, minimum width=3.2cm, minimum height=1.1cm] (spark) at (5.4,1.1) {Spark batch +\\Structured Streaming};
\node[f, minimum width=3.2cm, minimum height=1.1cm] (dash) at (5.4,-1.5) {Dashboard\\(bản đồ, cảnh báo)};
\foreach \s in {s1,s2,s3,s4,s5} { \draw[arr] (\s) -- (raw); }
\draw[arr] (raw) -- (ms);
\draw[arr] (ms) -- (spark);
\draw[arr] (ms) -- (dash);
\end{tikzpicture}%
}
\end{center}
\end{frame}

% Frame 8
\begin{frame}{Nhiệt độ bị đo trùng bởi cả 5 nguồn -- Survivorship chọn trạm làm ground truth trước}
\centering\small
\begin{tabular}{@{}lcccl@{}}
\toprule
Hạng ưu tiên & Nguồn & Bản chất & \multicolumn{2}{l}{Lý do xếp hạng} \\
\midrule
1 (cao nhất) & ISD / GHCN & Đo trực tiếp tại trạm & \multicolumn{2}{l}{Có QC flag, ground truth thật} \\
2 & ERA5 & Đồng hoá quan trắc + mô hình & \multicolumn{2}{l}{Độ phân giải cao nhất nhóm mô hình} \\
3 & Open-Meteo & Nội suy từ ERA5 & \multicolumn{2}{l}{Phái sinh từ hạng 2} \\
4 (thấp nhất) & NASA POWER & Vệ tinh/MERRA-2 & \multicolumn{2}{l}{Độ phân giải thô nhất (0.5\textdegree)} \\
\bottomrule
\end{tabular}
\vspace{0.35cm}

\small \textbf{Cảnh báo tích hợp sai}: \texttt{pressure\_surface} (POWER) và \texttt{pressure\_msl}
(Open-Meteo/ISD) là \textbf{2 đại lượng vật lý khác nhau} -- không được gộp trực tiếp, phải giữ
2 cột riêng trong Mediated Schema.
\end{frame}

%=============================================================
\section{Trạm quan trắc \& dung lượng}
%=============================================================
\hypertarget{sec4}{}
\hustsectionpage{4. Trạm quan trắc \&\\dung lượng dữ liệu}

% Frame 9
\begin{frame}{56 trạm GHCN Việt Nam và 246 trạm ISD toàn Đông Nam Á đã được lọc thật, có toạ độ}
\begin{center}
\bignum{56}{trạm GHCN-Daily\\toàn quốc Việt Nam}{hustblue}{3.6cm}
\bignum{10}{trạm ISD đang hoạt động\\(cập nhật đến 24/8/2025)}{hustblue}{3.6cm}
\bignum{18}{trạm lõi ưu tiên\\Tuần 1--2 (toàn ĐNA)}{goodgreen}{3.6cm}
\end{center}
\vspace{0.3cm}
\centering\scriptsize
\begin{tabular}{@{}lcccccccc@{}}
\toprule
Nước & VN & Indonesia & Thái Lan & Philippines & Malaysia & Lào & Myanmar & Campuchia/Brunei \\
\midrule
Số trạm GHCN & 56 & 104 & 54 & 34 & 16 & 15 & 5 & 1/1 \\
\bottomrule
\end{tabular}

\vspace{0.2cm}
\small \textcolor{badred}{Phát hiện khi lọc}: mã quốc gia Việt Nam trong ISD không nhất quán
(\texttt{VM} và \texttt{VN}); mã Brunei là \texttt{BX}, không phải \texttt{BS} (Bahamas).
\end{frame}

% Frame 10
\begin{frame}{Dung lượng dữ liệu trạm thực đo rất nhẹ -- ERA5 mới là phần quyết định tiêu chí Volume}
\centering\small
\begin{tabular}{@{}p{4.6cm}p{7.4cm}@{}}
\toprule
Nguồn (phạm vi đề xuất) & Dung lượng thật/tính toán \\
\midrule
GHCN 16 trạm VN, toàn lịch sử & \textbf{25,85 MB} (đã HEAD-check thật, TB 1,65 MB/trạm) \\
ISD 10 trạm VN, 1 năm (2024) & \textbf{42,94 MB} (đã HEAD-check thật, TB 4,29 MB/trạm) \\
ISD 10 trạm VN, 5 năm (ước) & $\approx$ 215 MB \\
NASA POWER + Open-Meteo & Nhẹ ($<$100 MB, truy vấn theo điểm) \\
\textbf{ERA5 (ARCO), vùng ĐNA, 8 biến, 5 năm, hourly} & \textbf{$\approx$ 34,8 GB} (tính từ chunk shape thật: 1440$\times$721, float32) \\
\bottomrule
\end{tabular}
\vspace{0.3cm}

\small Tổng ước tính \textbf{$\approx$ 35--36 GB} -- vượt khả năng xử lý bằng công cụ đơn máy
thông thường (đúng tiêu chí Volume của Big Data) nhưng vẫn khả thi để tải/lưu trong 5 tuần với
nhóm 3 người, thực tế hơn nhiều so với ước tính ban đầu 1,5--2 TB.
\end{frame}

%=============================================================
\section{Kế hoạch \& đội hình}
%=============================================================
\hypertarget{sec5}{}
\hustsectionpage{5. Kế hoạch \&\\đội hình}

% Frame 11
\begin{frame}{Nhóm 3 người thật -- một thành viên kiêm nhiệm cả Spark/PINN và Dashboard/PM}
\centering\small
\begin{tabular}{@{}p{3.0cm}p{3.6cm}p{5.6cm}@{}}
\toprule
Thành viên & Vai trò & Trọng tâm \\
\midrule
Ngô Văn Quang & TV1 -- Ingestion Lead & Thu thập \& làm sạch đa nguồn \\
Nguyễn Cao Lương & TV2 -- Schema Lead & Mediated Schema, Record Linkage \\
Trương Tuấn Nghĩa & TV3 + TV4 (kiêm nhiệm) & Spark/PINN \textbf{+} Dashboard/PM \\
\bottomrule
\end{tabular}
\vspace{0.35cm}

\small \textbf{Khuyến nghị giảm tải cho vai trò kiêm nhiệm}: coi PINN là điểm cộng có thể bỏ (đã
có phương án lùi về MLlib hồi quy); ưu tiên tuyệt đối Spark batch/streaming + Dashboard hoạt động
đúng trước khi đầu tư vào mô hình vật lý.
\end{frame}

% Frame 12
\begin{frame}{Timeline bám sát đúng 5 mốc chính thức của học phần, không có mốc tự đặt thêm}
\begin{center}
\resizebox{0.92\textwidth}{!}{%
\begin{tikzpicture}[remember picture,
  m/.style={rectangle, rounded corners, draw=hustblue, thick, fill=hustblue!8, minimum width=2.6cm, minimum height=1.1cm, align=center, font=\small},
  arr/.style={-{Latex[length=3mm]}, thick, hustblue}]
\node[m, fill=goodgreen!15, draw=goodgreen] (m1) at (-9,0) {26/9\\Chốt nhóm\\\textcolor{goodgreen}{\checkmark{} Đã qua}};
\node[m, fill=warnorange!15, draw=warnorange] (m2) at (-4.5,0) {5/10\\Nộp đề cương\\\textbf{Sắp tới}};
\node[m] (m3) at (0,0) {20/10\\Báo cáo tiến độ};
\node[m] (m4) at (4.5,0) {27/10\\Test};
\node[m, fill=hustblue!15] (m5) at (9,0) {8/11\\Báo cáo cuối kỳ};
\draw[arr] (m1) -- (m2); \draw[arr] (m2) -- (m3); \draw[arr] (m3) -- (m4); \draw[arr] (m4) -- (m5);
\end{tikzpicture}%
}
\end{center}
\vspace{0.4cm}
\centering\small
160 tác vụ chi tiết (4 tuần $\times$ 4 vai trò $\times$ 10 tác vụ/tuần), mỗi tác vụ có tiêu chí
kiểm thử (Definition of Done) -- xem \texttt{docs/03-ke-hoach-160-tasks.md}.
\end{frame}

% Frame 12b -- multi-track pipeline timeline, converges at each official milestone
\begin{frame}{Ba luồng công việc chạy song song suốt 4 tuần và hội tụ đúng tại từng mốc chính thức}
\vspace{-0.5cm}
\begin{center}
\resizebox{0.98\textwidth}{!}{%
\begin{tikzpicture}[remember picture,
  seg/.style={rectangle, draw=hustblue, thick, align=center, font=\scriptsize, minimum height=0.9cm, text width=1.9cm},
  lane/.style={font=\bfseries\small, color=hustblue},
  mile/.style={rectangle, rounded corners, draw=warnorange, thick, fill=warnorange!15, minimum height=0.7cm, align=center, font=\scriptsize},
  conv/.style={-{Latex[length=2.2mm]}, thick, goodgreen}]

% ---- lane labels ----
\node[lane, anchor=east] at (-0.3,3.4) {TV1};
\node[lane, anchor=east] at (-0.3,2.2) {TV2};
\node[lane, anchor=east] at (-0.3,1.0) {TV3+4};

% ---- TV1: Ingestion ----
\node[seg, minimum width=2.4cm, fill=hustblue!10] at (1.2,3.4) {Đăng ký API,\\crawl mẫu};
\node[seg, minimum width=3.8cm, fill=hustblue!18] at (4.3,3.4) {Batch crawl\\toàn bộ nguồn};
\node[seg, minimum width=2.0cm, fill=hustblue!10] at (7.2,3.4) {Veracity,\\outlier, Docker};
\node[seg, minimum width=3.2cm, fill=hustblue!18] at (9.8,3.4) {Data freeze,\\README};

% ---- TV2: Schema ----
\node[seg, minimum width=2.4cm, fill=goodgreen!10] at (1.2,2.2) {Thiết kế\\Mediated Schema};
\node[seg, minimum width=3.8cm, fill=goodgreen!18] at (4.3,2.2) {Warehouse,\\Kriging interp.};
\node[seg, minimum width=2.0cm, fill=goodgreen!10] at (7.2,2.2) {Record Linkage\\đầy đủ};
\node[seg, minimum width=3.2cm, fill=goodgreen!18] at (9.8,2.2) {Freeze schema,\\Data Dictionary};

% ---- TV3+4: Spark/PINN + Dashboard/PM ----
\node[seg, minimum width=2.4cm, fill=warnorange!10] at (1.2,1.0) {Setup Spark,\\roadmap};
\node[seg, minimum width=3.8cm, fill=warnorange!18] at (4.3,1.0) {Streaming,\\Dashboard khung};
\node[seg, minimum width=2.0cm, fill=warnorange!10] at (7.2,1.0) {PINN,\\Dashboard nâng cấp};
\node[seg, minimum width=3.2cm, fill=warnorange!18] at (9.8,1.0) {Inference,\\hoàn thiện Dashboard};

% ---- milestone verticals ----
\foreach \x/\lbl in {2.4/{5/10\\Đề cương}, 6.2/{20/10\\Tiến độ}, 8.2/{27/10\\Test}} {
  \draw[dashed, warnorange, thick] (\x,0.45) -- (\x,4.05);
  \node[mile] at (\x,4.35) {\lbl};
}
\node[mile, fill=goodgreen!20, draw=goodgreen] at (0,4.35) {26/9\\\checkmark{} Chốt nhóm};
\draw[dashed, goodgreen, thick] (0,0.45) -- (0,4.05);

% ---- convergence into final milestone ----
\draw[conv] (11.4,3.4) -- (12.5,2.2);
\draw[conv] (11.4,2.2) -- (12.5,2.2);
\draw[conv] (11.4,1.0) -- (12.5,2.2);
\node[rectangle, rounded corners, draw=goodgreen, very thick, fill=goodgreen!15,
      minimum width=2.1cm, minimum height=1.6cm, align=center, font=\bfseries\scriptsize]
      at (13.7,2.2) {8/11\\BÁO CÁO\\CUỐI KỲ};
\end{tikzpicture}%
}
\end{center}
\vspace{0.2cm}
\centering\scriptsize
Mỗi luồng độc lập theo tuần nhưng \textbf{cùng đứng trước mỗi mốc chính thức} -- không luồng nào
được phép trễ hơn mốc chung, vì Spark (TV3+4) cần dữ liệu đã chuẩn hoá của TV1/TV2 làm đầu vào.
\end{frame}

%=============================================================
\section{Rủi ro \& kết luận}
%=============================================================
\hypertarget{sec6}{}
\hustsectionpage{6. Rủi ro \&\\kết luận}

% Frame 13
\begin{frame}{Ba rủi ro lớn nhất đều đã có phương án dự phòng cụ thể, không phải giả định suông}
\centering\footnotesize
\begin{tabular}{@{}p{3.3cm}p{3.3cm}p{5.4cm}@{}}
\toprule
Rủi ro & Mức ảnh hưởng & Phương án dự phòng \\
\midrule
Đăng ký CDS chậm & Trung bình & Dùng ARCO-ERA5 (Google Cloud, miễn phí, không đăng ký) song song \\
Quá tải vai trò kiêm nhiệm & Cao & Bỏ PINN nếu cần, giữ Spark + Dashboard làm lõi bắt buộc \\
Độ phủ trạm lệch giữa các nước & Thấp--Trung bình & Nêu rõ giới hạn (Campuchia/Brunei chỉ 1 trạm/nước) trong Veracity report \\
\bottomrule
\end{tabular}
\end{frame}

% Frame 14
\begin{frame}{Đề tài đã có đủ nguồn dữ liệu thật, kiến trúc cụ thể, và kế hoạch khả thi cho 3 người}
\begin{center}
\bignum{4/4}{nguồn lõi sẵn sàng dùng ngay\\hôm nay, không chờ thủ tục}{goodgreen}{4.2cm}
\bignum{18}{trạm quan trắc thật\\đã có toạ độ, ID cụ thể}{goodgreen}{4.2cm}
\bignum{35 GB}{tổng dữ liệu cần thiết\\ước tính có căn cứ, khả thi}{goodgreen}{4.2cm}
\end{center}
\vspace{0.4cm}
\centering\small
Toàn bộ tài liệu chi tiết (đề xuất, nguồn dữ liệu, kế hoạch, giải pháp tích hợp, danh sách trạm)
nằm tại repository GitHub riêng của nhóm -- sẵn sàng cho buổi nộp đề cương 5/10.
\end{frame}

\hustthankyou

\end{document}
